\section*{Capítulo \ref{ch:b1:stratigraphic-scale} --- La escala estratigráfica}

\begin{solution}{exo:b1:stratigraphic-scale:1}
(a)~Cronoestratigráfica: un \omterm{def:b1:stratigraphic-scale:chronostratigraphy}{sistema}. (b)~Litoestratigráfica: una \omterm{def:b1:stratigraphic-scale:lithostratigraphy}{formación}.
(c)~Geocronológica: una época (su roca es la \omterm{def:b1:stratigraphic-scale:chronostratigraphy}{Serie} Cretácica Inferior).
(d)~Bioestratigráfica: una \omterm{def:b1:stratigraphic-scale:biozone}{biozona}. (e)~Cronoestratigráfica: un \omterm{def:b1:stratigraphic-scale:chronostratigraphy}{piso}.
(f)~Geocronológica: una época. (g)~Litoestratigráfica: un \omterm{def:b1:stratigraphic-scale:lithostratigraphy}{miembro}.
\end{solution}

\begin{solution}{exo:b1:stratigraphic-scale:2}
Paleozoico $538.8 - 251.902 = \qty{286.9}{\Myr}$, el \qty{53.2}{\%} del
Fanerozoico; Mesozoico $251.902 - 66.00 = \qty{185.9}{\Myr}$, el \qty{34.5}{\%};
Cenozoico \qty{66.0}{\Myr}, el \qty{12.2}{\%}. El Fanerozoico es
$538.8/4567 = \qty{11.8}{\%}$ del tiempo desde la base del Hadeico.
\end{solution}

\begin{solution}{exo:b1:stratigraphic-scale:3}
A, B, C y D se depositaron en ese orden (superposición), horizontales (horizontalidad original), así que su inclinación es posterior. El dique E corta de A a D, así que es más joven que D
(relaciones de corte); si fue anterior o posterior a la inclinación no puede decirse por la descripción. La superficie de erosión corta D y E, así que es más joven que ambos
(relaciones de corte). F y G se depositaron después sobre ella (superposición). La falla desplaza de A a F pero no a G: se movió después de F y antes de G
(relaciones de corte). Orden: A, B, C, D; inclinación e intrusión de E (orden no decidido); erosión; F; falla; G.
\end{solution}

\begin{solution}{exo:b1:stratigraphic-scale:4}
$r = 30/(85.40 - 84.60) = \qty{37.5}{m}$ por millón de años. La primera aparición
yace $12/37.5 = \qty{0.32}{\Myr}$ después de la ceniza inferior: $85.40 - 0.32 =
\qty{85.08}{\Ma}$.
\end{solution}

\begin{solution}{exo:b1:stratigraphic-scale:5}
(a)~Cuatro barras verticales frente a la altura. (b)~B y C aparecen juntos entre \qtyrange{26}{50}{m}: una \omterm{def:b1:stratigraphic-scale:biozone}{zona de rango} concurrente de \qty{24}{m} de espesor. La \omterm{def:b1:stratigraphic-scale:biozone}{zona de intervalo}
desde la primera aparición de B hasta la de D va de \qtyrange{12}{41}{m}, \qty{29}{m}. (c)~En la segunda sección pueden situarse la primera aparición de B (\qty{3}{m}) y la de C (\qty{9}{m}), así que se reconocen la base de la \omterm{def:b1:stratigraphic-scale:biozone}{zona de intervalo} de (b) y
la base de la \omterm{def:b1:stratigraphic-scale:biozone}{zona de rango} concurrente; su techo (la primera aparición de D) y el techo del solapamiento tal como se ven en la sección~1 no pueden comprobarse, pues D está ausente. Las primeras apariciones de B y C ($\qty{12}{m} \leftrightarrow
\qty{3}{m}$, $\qty{26}{m} \leftrightarrow \qty{9}{m}$) son las líneas más seguras; las últimas apariciones (de B, a \qty{50}{m}, y de B y C, a \qty{15}{m}) dependen más de la preservación.
\end{solution}

\begin{solution}{exo:b1:stratigraphic-scale:6}
(a)~$60/20 = \qty{3}{\Myr}$: el techo de la arenisca es tres millones de años más joven en el pozo tierra adentro. (b)~Un millón de años de lodo a \qty{30}{m} por millón de años: la ceniza yace
\qty{30}{m} por encima del techo de la arenisca. (c)~Cuando cayó la ceniza, la línea de costa se había desplazado \qty{20}{km} tierra adentro del pozo más marino, así que seguía \qty{40}{km} mar adentro del pozo tierra adentro, que estaba sobre tierra firme. La ceniza cayó sobre la superficie de tierra allí, por debajo de la arenisca que el mar depositó dos millones de años después ---si no fue arrastrada antes, como puede serlo una capa suelta en tierra. La misma capa de ceniza está por encima del límite de \omterm{def:b1:stratigraphic-scale:lithostratigraphy}{formación} en un pozo y por debajo en el otro: el límite cruza el tiempo.
\end{solution}

\begin{solution}{exo:b1:stratigraphic-scale:7}
(a)~Lochkoviense $1.36/419.62 = \qty{0.32}{\%}$; induense $0.024/251.902 =
\qty{0.0095}{\%}$, unas 34 veces menor. La base del induense yace entre capas de ceniza datadas con precisión en la propia sección del GSSP; una incertidumbre grande sugiere un límite cuya edad debe transportarse desde niveles datados más lejos.
(b)~$251.941 - 251.880 = \qty{0.061}{\Myr} = 61$ mil años, con la incertidumbre $\sqrt{37^2 + 31^2} = 48$ mil años: los
\num[uncertainty-mode=separate]{60 \pm 48} mil años publicados.
\end{solution}

\begin{solution}{exo:b1:stratigraphic-scale:8}
(a)~$m = 3000 \times \frac43\pi\,(5000)^3 = \qty{1.57e15}{kg}$; el polvo, $60m = \qty{9.42e16}{kg}$. (b)~La superficie de la Tierra es $4\pi(6.378\times10^6)^2
= \qty{5.11e14}{m^2}$: \qty{184}{kg} por metro cuadrado. (c)~$184/2500 =
\qty{0.074}{m}$, unos \qty{7}{cm}. (d)~El espesor varía como el cubo del diámetro: \qty{1.6}{cm} para \qty{6}{km}, \qty{20}{cm} para \qty{14}{km}. El orden de magnitud, de centímetros a decímetros, es el de la arcilla del límite,
\qty{50}{cm} de espesor en El Kef; la comparación solo puede ser aproximada, pues el polvo no pudo caer de manera uniforme y una arcilla también recoge el sedimento ordinario de su tiempo.
\end{solution}

\begin{solution}{exo:b1:stratigraphic-scale:9}
(a)~$\num{60e3} \times \num{2e-6} = \qty{0.12}{m}$, \qty{12}{cm}; para 12 y 108 mil años, \qty{2.4}{cm} y \qty{21.6}{cm}. (b)~$\num{32e3} \times
\num{10e-6} = \qty{0.32}{m}$. (c)~Una extinción ocupa de centímetros a decímetros de roca: las muestras deben tomarse centímetro a centímetro, y una sola muestra de \qty{10}{cm} de espesor en la sección condensada promedia cincuenta mil años ---casi todo el suceso.
\end{solution}

\begin{solution}{exo:b1:stratigraphic-scale:10}
(a)~$u = R/s = 10 \times 2000 = \qty{20000}{m}$ por millón de años,
\qty{20}{km} por millón de años; en \qty{30}{km}, $\Delta t = \qty{1.5}{\Myr}$.
(b)~La llanura es ahora la superficie $z = s x + a t$. La línea de costa
$s x + a t = z_0 + R t$, así que $x(t) = (z_0 + (R - a)t)/s$ y
$\mathrm dx/\mathrm dt = (R - a)/s$, positiva (hacia tierra) mientras $a < R$. Cuando
$a > R$ la velocidad es negativa: la línea de costa se desplaza mar adentro, una regresión, y la arenisca de la costa es entonces más joven mar adentro. (c)~Un acantilado de \qty{200}{m} de longitud abarca, a \qty{20}{km} por millón de años, $0.2/20 = \qty{0.01}{\Myr}$ --- diez mil años de diacronismo, muy por debajo de lo que resuelven los fósiles o la mayoría de las dataciones; la roca a uno y otro lado del contacto es la misma en toda su extensión.
\end{solution}

\begin{solution}{exo:b1:stratigraphic-scale:11}
(a)~$f = (25 - 10)/(50 - 10) = 0.375$ y $t = 100.50 - 0.375 \times 0.80 =
\qty{100.20}{\Ma}$, con
\[
  \sigma_t = \sqrt{(0.625 \times 0.05)^2 + (0.375 \times 0.03)^2} = \qty{0.033}{\Myr}.
\] (b)~$g(f) = (1 - f)^2\sigma_1^2 + f^2\sigma_2^2$
tiene $g'(f) = -2(1 - f)\sigma_1^2 + 2f\sigma_2^2$, cero en $f^* =
\sigma_1^2/(\sigma_1^2 + \sigma_2^2) = 0.0025/0.0034 = 0.735$, más cercano a la capa mejor
datada; allí $\sigma_t = \sigma_1\sigma_2/\sqrt{\sigma_1^2 + \sigma_2^2} =
\qty{0.026}{\Myr}$, menor que cualquiera de las dos. (c)~$g$ es una parábola que se abre hacia arriba, así que en $[0, 1]$ es mayor en un extremo: $g(0) = \sigma_1^2$, $g(1) = \sigma_2^2$.
Por tanto $\sigma_t \le \max(\sigma_1, \sigma_2)$ entre las capas.
\end{solution}

\begin{solution}{exo:b1:stratigraphic-scale:12}
(a)~El límite está a mitad de profundidad: $f = 0.5$, \qty{29.50}{\Ma}. (b)~Con
$\varphi_0\lambda = 0.63 \times 1960 = \qty{1234.8}{m}$, el espesor sólido de
\qtyrange{600}{1000}{m} es $400 - 1234.8\,(e^{-600/1960} - e^{-1000/1960}) =
\qty{232.2}{m}$, y de \qtyrange{200}{600}{m}, \qty{194.2}{m}; en total
\qty{426.3}{m}. La fracción de los granos depositados entre la ceniza inferior y el
límite es $232.2/426.3 = 0.545$, y la edad es $30.0 - 0.545 =
\qty{29.46}{\Ma}$. (c)~La segunda estimación es \num{45} mil años más joven.
Es la mejor: la mitad más profunda ha perdido más agua, así que contiene más granos, y más tiempo, por metro que la mitad más somera (porosidad media
\num{0.42} frente a \num{0.51}).
\end{solution}

\subsection*{Problema: correlación de tres secciones a través de una cuenca}

\begin{solution}{pb:b1:stratigraphic-scale:1}
\begin{enumerate}
  \item Tres columnas, digamos a \qty{1}{cm} por \qty{10}{m}, unidas por las capas de ceniza y por los límites de zona; la cuña de la Arenisca de Kell se adelgaza de \qty{45}{m} en A a \qty{6}{m} en C.
  \item Zona~2: A $65.8 - 36.6 = \qty{29.2}{m}$, B \qty{20.8}{m}, C \qty{10.0}{m}. Zona~3 registrada: A \qty{14.2}{m}, B \qty{14.0}{m}, C \qty{7.2}{m}.
  \item La zona~1 comienza en la primera aparición de \emph{V.~primum}, que se encuentra desde la base de cada sección: su primera aparición yace por debajo de todas ellas.
  \item Solo vive en la arenisca: sigue a la facies, y la facies es de edades distintas de un lugar a otro. En A aparece en las zonas~1 y~2, en B y C solo en la zona~1.
  \item En la zona~2 en las tres (A \qty{49.0}{m}, B \qty{40.0}{m}, C \qty{19.0}{m}). Es coherente de sección en sección aquí, y así podría subdividir la zona~2, pero una última aparición depende de la preservación; debería comprobarse antes en más secciones.
  \item Las líneas unen \qty{36.6}{m}, \qty{31.2}{m} y \qty{14.8}{m}, y después \qty{65.8}{m}, \qty{52.0}{m} y \qty{24.8}{m}. La zona~2 se adelgaza en la proporción $29.2 : 20.8 : 10.0$, aproximadamente $2.9 : 2.1 : 1$: si la zona duró el mismo tiempo en todas partes, el sedimento se acumuló unas tres veces más deprisa en A que en C.
  \item A: zona~2 ($45.0 > 36.6$); B: zona~1 ($20.0 < 31.2$); C: zona~1 ($6.0 < 14.8$).
  \item No: está en la zona~2 en A y en la zona~1 en B y C, así que es más joven hacia el oeste. La lutita marina se extendió hacia el oeste sobre la arena de la costa: la línea de costa se desplazó hacia el oeste, tierra adentro, en una transgresión.
  \item Una ceniza cae sobre toda la cuenca en un instante geológico y se asienta sobre cualquier sedimento; el techo de la arenisca sigue a la línea de costa en movimiento y es \omterm{def:b1:stratigraphic-scale:diachronous}{diacrónico}.
  \item Continuidad lateral (la capa continúa a través de la cuenca) y superposición (la ceniza~1 bajo la ceniza~2 en todas las secciones, lo que comprueba la correlación).
  \item A $(71.0 - 5.0)/1.20 = \qty{55.0}{m}$ por millón de años; B $48.0/1.20 = \qty{40.0}{}$; C $24.0/1.20 = \qty{20.0}{}$.
  \item A es la más cercana a la tierra, de donde procedía el aporte de sedimento; llegó menos hasta las aguas profundas del este.
  \item $0.1/55 = \num{1.8}$ mil años en A, $\num{2.5}$ en B, $\num{5.0}$ en C: A resuelve el tiempo con mayor finura.
  \item A $152.40 - 40.0/55.0 = \qty{151.67}{\Ma}$; B $152.40 - 12.0/40.0 =
        \qty{152.10}{\Ma}$; C $152.40 - 3.0/20.0 = \qty{152.25}{\Ma}$.
  \item C a B: $20/(152.25 - 152.10) = \qty{133}{km}$ por millón de años; B a A:
        $25/(152.10 - 151.67) = \qty{58.5}{km}$ por millón de años. Se frenó.
  \item A $152.40 - 60.8/55.0 = \qty{151.29}{\Ma}$; B $152.40 - 44.0/40.0 =
        \qty{151.30}{\Ma}$; C $152.40 - 21.8/20.0 = \qty{151.31}{\Ma}$.
  \item Abarcan \num{16} mil años, menos que la incertidumbre de cada una (unos \num{28} mil años, pregunta~19): la primera aparición es sincrónica en toda la cuenca dentro de la resolución de las dataciones, y eso la convierte en un buen índice.
  \item A $152.40 - 31.6/55.0 = \qty{151.826}{\Ma}$; B $152.40 - 23.2/40.0 =
        \qty{151.820}{\Ma}$; C $152.40 - 11.8/20.0 = \qty{151.810}{\Ma}$.
  \item $f = 0.479$, $0.483$, $0.492$; $\sigma_t = \sqrt{(1 - f)^2 + f^2} \times
        0.04 = \qty{0.028}{\Myr}$ en cada una, \num{28} mil años.
  \item Media \qty{151.819}{\Ma}; dispersión $151.826 - 151.810 = \num{16}$ mil años, menor que la incertidumbre de cada edad: las tres concuerdan.
  \item La interpolación supone una tasa única desde la ceniza~1 hasta la ceniza~2. En A ese intervalo contiene arenisca de costa hasta \qty{45}{m} y lutita por encima; la arena y el lodo no tienen por qué acumularse a la misma tasa, y se compactan de forma distinta (\cref{ch:b1:diagenesis-basins}). Si la arena se acumuló más deprisa, la parte arenosa contiene menos tiempo del que sugiere su espesor, y la primera aparición es más antigua de lo calculado: la edad de A es, en efecto, la más antigua de las tres.
  \item B: la primera aparición yace en lutita marina, \qty{11}{m} por encima del techo de la arenisca, sin cambio de facies cerca del nivel, y la tasa da \num{2.5} mil años por \qty{10}{cm}. A resuelve mejor el tiempo, pero sitúa el nivel en la arenisca de costa, \qty{8.4}{m} por debajo de un cambio de facies; C está condensada, \num{5} mil años por \qty{10}{cm}. C es la peor para la resolución, y A para las facies.
  \item A partir de las edades medias de las dos primeras apariciones, $151.819 - 151.302 =
        \qty{0.52}{\Myr}$: unos medio millón de años.
  \item El límite es \qty[uncertainty-mode=separate]{151.82(3)}{\Ma}.
\end{enumerate}
\end{solution}
